\subsection{S4: conservative development successor}
After S3 was closed, a separately declared S4 development study reused the already inspected 2023/2024 public windows. It is not a fresh generalization test. Six recent-base fits (three seeds per window) generated 24 prediction cells. The base used the most recent 96 pre-validation dates, divided into 48/12/12/24 fitting, checkpoint-selection, calibration and residual dates with strict maturity purges, and five training epochs. All arms shared this base and mean-only correction. The distance intervention multiplies keys by $\sqrt{d}$, exactly converting the inherited mean squared distance to squared Euclidean distance. Temperature is fixed at one.

The four arms are memory off, original distance, Euclidean distance and matched random. Nonzero weights from $\{0.25,0.5,1\}$ were eligible only if each chronological half of validation improved NLL over zero correction by more than 0.01 nat. Every non-off arm in all six seed/window cells selected zero. Consequently all four arms scored the same NLL: -1.332779 in 2023 and -1.411278 in 2024. No memory improvement was demonstrated. Because the gate is zero, deployed equality does not establish that distance scaling itself is ineffective.

All 247 pre-fit tests passed. The campaign took 64.76 seconds and reached 401.84 MiB sampled worker RSS under the explicitly renewed single-worker slow-mode exception. An independent arithmetic audit reconstructed 3,168 rows to $2.0\times10^{-15}$ and verified 81 frozen source/data identities. These are development and engineering observations, not current market forecasts, protected confirmation or independent retraining. Source: \path{outputs/frem_s4_20261012_attempt1/completion.json}; audit: \path{research/frem_s4_20261012/audit.json}. Completed S2/S3 results and all predecessor artifacts remain preserved.
